\begin{table}[htbp]
\centering
\small
\caption{Шесть ревизий манифеста допуска и причины перехода}
\label{tab:s12}
\begin{tabular}{>{\RaggedRight\hspace{0pt}\arraybackslash}p{0.460\linewidth}>{\RaggedRight\hspace{0pt}\arraybackslash}p{0.460\linewidth}}
\toprule
Ревизия & Что заставило перейти к следующей \\
\midrule
1 & правка clf\_run.py по амендменту о двух схемах inner CV \\
2 & процедура 1 не создавала дочерний манифест — цепочка была разорвана \\
3 & вызов pick\_c в P2b остался с прежним числом возвращаемых значений и останавливал прогон \\
4 & отчёты fairness и разбора ошибок называли себя серией v1 и ссылались на файл сверки v1, хотя сверялись с v2 \\
5 & разбор ошибок переобучал модель со своим GroupKFold вместо зарегистрированной схемы \\
6 & действующая \\
\bottomrule
\end{tabular}
\begin{minipage}{\linewidth}\footnotesize\vspace{2pt}
\textit{Примечание.} Единица анализа — ревизия манифеста. Знаменатель: шесть ревизий за 29 июля 2026. Данные: манифесты допуска серии v2. Таблица описывает историю прогонов; чисел в ней нет. Шкала: шкалы нет. Сокращения: inner CV — вложенная кросс-валидация; GroupKFold — разбиение с непересекающимися группами. Пропуски: ошибочный манифест не удаляется, а помечается invalidated с причиной, поэтому строк ровно столько, сколько было ревизий. Поправка на множественные сравнения не применялась.
\end{minipage}
\end{table}
